\chapter{Dérivés sur matières premières et énergie}\label{ch:rc:commodity-and-energy-derivatives}

Pendant la deuxième semaine de février 2021, une vague de froid gela les puits de gaz à travers le Texas et
les grandes plaines du Sud. Le gaz au comptant du hub de Waha, dans l'ouest du Texas, atteignit 206.19 dollars par
million de British thermal units le 16 février~; à un hub de l'Oklahoma, le prix moyen du 17 février fut
de 1\,192 dollars, contre moins de 3 dollars la première semaine du mois,
selon l'US Energy Information Administration. Henry Hub, l'indice de référence, passa
de 2.88 dollars le premier jour de février à 23.86 le 17, puis redescendit sous 3 dollars le 23. Pendant
une semaine, la valeur du gaz dans le sous-sol dépendait de la vitesse à laquelle on pouvait le retirer~: un
propriétaire d'une cavité saline pouvant se vider en quelques jours pouvait vendre au sommet, celui
d'un réservoir lent ne le pouvait pas. Ce chapitre évalue cette flexibilité. Le livre~3 a décrit
les courbes à terme, l'économie du stockage et les contrats structurés des marchés de matières premières~;
on leur donne ici des modèles~: une courbe à terme qui se déplace selon deux facteurs, des options sur l'écart
entre deux prix, et l'évaluation des droits de stockage et de swing par optimisation sous
incertitude.

\section{Modèles factoriels de la courbe à terme}

Une matière première n'a pas un prix mais une courbe à terme~: chaque mois de livraison se négocie, avec
une forme saisonnière et un rendement d'opportunité (convenience yield) (livre~3, chapitre~10). Un modèle doit faire bouger
toute la courbe, et il doit faire en sorte que les contrats à courte échéance bougent davantage que ceux à longue échéance.

\begin{definition}[modèle de Schwartz--Smith]\label{def:rc:commodity-and-energy-derivatives:ss}
Le \emph{modèle de Schwartz--Smith}\index{modèle de Schwartz--Smith} écrit le logarithme du prix au comptant
comme la somme d'un écart de court terme $\chi_t$, un processus d'Ornstein--Uhlenbeck revenant vers
zéro à la vitesse $\kappa$ avec une volatilité $\sigma_s$, et d'un niveau de long terme $\xi_t$, un mouvement brownien
de volatilité $\sigma_l$, les deux étant corrélés par $\rho$. Sa courbe à terme évolue selon
\[
\frac{dF(t,T)}{F(t,T)} = \sigma_se^{-\kappa(T-t)}\,dW^1_t+\sigma_l\,dW^2_t,\qquad d\langle W^1,W^2\rangle = \rho\,dt,
\]
de sorte que $\ln F(t,T) = \ln F(0,T)+e^{-\kappa(T-t)}\chi_t+\xi_t-\tfrac12V(t,T)$, où $V$ désigne la variance.
La saisonnalité réside dans la courbe initiale $F(0,T)$, que le modèle ajuste exactement.
\end{definition}

Les deux facteurs sont gaussiens et se simulent exactement, et chaque terme est lognormal~: une option
sur le terme $F(\cdot,T)$ arrivant à échéance en $t$ se valorise par la formule de Black avec la
volatilité implicite $\sqrt{V(t,T)/t}$.

\begin{example}[calibration]\label{ex:rc:commodity-and-energy-derivatives:calib}
Une courbe illustrative de type Henry Hub va de 2.60 dollars pour avril à 3.50 pour janvier.
Des options arrivant à échéance une quinzaine de jours avant chaque livraison sont cotées à des volatilités implicites
allant de 62.0~\% (mai) à 43.0~\% (mars), données synthétiques. Avec $\rho = 30~\%$ fixé, un
ajustement de Levenberg--Marquardt donne $\kappa = 1.50$, $\sigma_s = 60.0~\%$ et $\sigma_l = 20.0~\%$.
\end{example}

\section{L'effet Samuelson et la volatilité à terme}

Dans le modèle, la volatilité instantanée d'un terme dont le délai avant livraison est $\tau$ vaut
\[
\sigma(\tau) = \sqrt{\sigma_s^2e^{-2\kappa\tau}+\sigma_l^2+2\rho\sigma_s\sigma_le^{-\kappa\tau}}:
\]
68.7~\% à la livraison, 27.2~\% à un an
de la livraison et 20.0~\% à cinq ans de la livraison, pour les paramètres calibrés. C'est l'effet Samuelson
du livre~3 traduit en propriété du modèle~: les chocs d'offre et de demande font bouger les contrats proches,
que porte le facteur de court terme, tandis que les contrats éloignés ne bougent qu'avec le
niveau de long terme (\cref{fig:rc:commodity-and-energy-derivatives:vol}). La corrélation entre deux
contrats décroît avec la distance entre leurs livraisons, et c'est ce qui donne une valeur aux \omterm{def:rc:commodity-and-energy-derivatives:spread}{options
sur spread}.

\begin{omfigure}
\begin{tikzpicture}
\begin{axis}[width=11.5cm,height=5.4cm,
  xlabel={délai avant livraison (années)},ylabel={volatilité (\%)},
  tick label style={font=\small},label style={font=\small},
  xmin=0,xmax=1,ymin=0,ymax=75,grid=major,grid style={gray!25},
  legend columns=2,legend style={font=\footnotesize,at={(0.5,-0.25)},anchor=north,draw=none,fill=none,/tikz/every even column/.append style={column sep=8pt}},
  table/col sep=comma]
\addplot[omDef,very thick] table[x=tau,y=vol]{figdata/rates-credit-risk/16-commodity-and-energy-derivatives/vol.csv};
\addplot[omThm,only marks,mark=*,mark size=1.8pt] table[x=t,y=vol]{figdata/rates-credit-risk/16-commodity-and-energy-derivatives/volquotes.csv};
\legend{volatilité instantanée à terme,cotations de volatilité implicite (à la date de livraison)}
\end{axis}
\end{tikzpicture}

\omcaption{Volatilité instantanée à terme du modèle calibré en fonction du délai avant
livraison, et les volatilités implicites synthétiques auxquelles il a été ajusté, tracées à l'instant de livraison de chaque
contrat. Une volatilité implicite moyenne la volatilité instantanée sur
la durée de vie de l'option, pour des délais avant livraison allant de tout l'horizon jusqu'à une quinzaine de jours, de sorte que
elle se situe au-dessus de la valeur de la courbe à son instant de livraison. Données~: le
tutoriel du chapitre.}\label{fig:rc:commodity-and-energy-derivatives:vol}
\end{omfigure}

\section{Options sur spread}

\begin{definition}[Option sur spread]\label{def:rc:commodity-and-energy-derivatives:spread}
Une \emph{option sur spread}\index{option sur spread} verse $\max(F_1-F_2-K,0)$ sur la différence
entre deux prix~: deux mois de livraison d'une même matière première (un calendar spread), un produit
et son intrant (un crack ou spark spread), ou une même matière première en deux lieux.
\end{definition}

Avec $K = 0$ et des prix lognormaux, l'option est une \omterm{def:rc:bermudans-and-callables:switch}{option d'échange}, évaluée exactement par
la formule de Margrabe avec la volatilité du ratio,
$\sigma^2 = \sigma_1^2+\sigma_2^2-2\rho\sigma_1\sigma_2$. Avec $K\ne0$, l'écart n'est pas lognormal et
il n'existe pas de forme fermée.

\begin{definition}[Approximation de Kirk]\label{def:rc:commodity-and-energy-derivatives:kirk}
L'\emph{approximation de Kirk}\index{approximation de Kirk} évalue un call sur spread en traitant
$F_2+K$ comme lognormal de volatilité $\sigma_2F_2/(F_2+K)$, puis en appliquant la formule de Margrabe~:
\[
C \approx P\bigl(F_1N(d_1)-(F_2+K)N(d_2)\bigr),\quad \sigma_K^2 = \sigma_1^2-2\rho\sigma_1\sigma_2b+\sigma_2^2b^2,\quad b = \frac{F_2}{F_2+K}.
\]
\end{definition}

\omcode{code/firm/energymodel/firm_energymodel.py}{102}{107}{Approximation de Kirk d'un call sur spread.}\label{lst:rc:commodity-and-energy-derivatives:kirk}

\begin{example}[Un calendar spread]\label{ex:rc:commodity-and-energy-derivatives:cal}
Un call sur le terme de janvier (3.50) moins le terme de juillet (2.74), de prix d'exercice 0.50, arrivant
à l'échéance de livraison de juillet~: le modèle donne des volatilités de 34.8~\% et 58.1~\% à cette date et une
corrélation de 0.970. L'\omterm{def:rc:commodity-and-energy-derivatives:kirk}{approximation de Kirk} l'évalue à 29.41 cents par MMBtu, contre
29.78 par la méthode de Monte-Carlo (erreur type 0.01), soit une erreur de 0.38 cent~; à strike nul, les deux donnent
environ 75.8 cents, la valeur de Margrabe. L'erreur est maximale près de la monnaie, là où l'approximation de la
loi de la somme joue le plus (\cref{fig:rc:commodity-and-energy-derivatives:kirk}).
\end{example}

\begin{omfigure}
\begin{tikzpicture}
\begin{axis}[width=11.5cm,height=5.2cm,
  xlabel={strike $K$ (\$/MMBtu)},ylabel={écart Kirk moins Monte Carlo (cents)},
  tick label style={font=\small},label style={font=\small},
  xmin=-0.5,xmax=1.5,ymin=-0.5,ymax=0.5,grid=major,grid style={gray!25},
  legend columns=2,legend style={font=\footnotesize,at={(0.5,-0.25)},anchor=north,draw=none,fill=none,/tikz/every even column/.append style={column sep=8pt}},
  table/col sep=comma]
\addplot[omDef,very thick,mark=*,mark size=1.2pt] table[x=k,y=err]{figdata/rates-credit-risk/16-commodity-and-energy-derivatives/kirk.csv};
\addplot[omExo,dashed] table[x=k,y=se]{figdata/rates-credit-risk/16-commodity-and-energy-derivatives/kirk.csv};
\addplot[omExo,dashed,forget plot] table[x=k,y expr=-\thisrow{se}]{figdata/rates-credit-risk/16-commodity-and-energy-derivatives/kirk.csv};
\legend{erreur de Kirk,$\pm2$ erreurs types de Monte Carlo}
\end{axis}
\end{tikzpicture}

\omcaption{Erreur de l'\omterm{def:rc:commodity-and-energy-derivatives:kirk}{approximation de Kirk} pour le call sur calendar spread janvier--juillet
face à 400\,000 trajectoires de Monte-Carlo, par strike. L'erreur est nulle à strike nul (la formule de Margrabe
est exacte) et change de signe autour de la monnaie~; elle reste inférieure à un demi-cent à chaque strike,
sur des options valant de zéro à 1.26~dollars. Données~: le tutoriel du chapitre.}\label{fig:rc:commodity-and-energy-derivatives:kirk}
\end{omfigure}

\section{Valorisation du stockage}

Une installation de stockage achète du gaz quand il est bon marché, le conserve et le vend quand il est cher,
dans les limites de capacité et de débits d'injection et de soutirage, et à un coût par unité
déplacée. Son propriétaire détient un portefeuille d'options sur calendar spread liées par l'inventaire.

\begin{definition}[Valeur intrinsèque du stockage]\label{def:rc:commodity-and-energy-derivatives:intrinsic}
La \emph{valeur intrinsèque du stockage}\index{valeur intrinsèque du stockage} d'une installation est la
valeur du meilleur calendrier d'injection et de soutirage sur la courbe à terme du jour, figée
par la négociation des contrats à terme~: une optimisation déterministe.
\end{definition}

\begin{definition}[Valeur extrinsèque, intrinsèque glissant]\label{def:rc:commodity-and-energy-derivatives:extrinsic}
La \emph{valeur extrinsèque}\index{valeur extrinsèque} d'une installation est sa valeur au-delà de
l'intrinsèque, issue du droit de modifier le calendrier au gré des prix. La stratégie
\emph{intrinsèque glissant}\index{intrinsèque glissant} (rolling intrinsic) en capte une partie~: on réoptimise le calendrier
sur chaque nouvelle courbe à terme et l'on couvre de nouveau l'écart, sans jamais perdre l'intrinsèque déjà
figé.
\end{definition}

\begin{method}[Stockage par Monte Carlo à moindres carrés]\label{met:rc:commodity-and-energy-derivatives:lsm}
On simule les prix spot pour chaque période. On remonte dans le temps à partir de la dernière période, la valeur
de chaque inventaire final étant connue sur chaque trajectoire. À chaque période et pour chaque niveau d'inventaire
atteignable après l'action, on régresse la valeur réalisée de la période suivante sur $1$, $S$ et $S^2$
pour estimer la valeur de continuation~; on choisit, trajectoire par trajectoire, l'action qui maximise le gain
immédiat plus la continuation estimée, et l'on reporte en arrière la valeur réalisée de ce choix. La
valeur est la moyenne au niveau d'inventaire initial.
\end{method}

\omcode{code/firm/energymodel/firm_energymodel.py}{169}{200}{Monte Carlo à moindres carrés pour le stockage~: régression des valeurs de continuation par inventaire final, et valeurs réalisées reportées en arrière.}\label{lst:rc:commodity-and-energy-derivatives:lsm}

\begin{example}[Une saison de stockage]\label{ex:rc:commodity-and-energy-derivatives:storage}
Une installation de 1 million de MMBtu (dix unités) démarre et termine l'année d'avril à mars vide~; elle
peut injecter deux unités par mois et en soutirer quatre, à 2~cents par MMBtu dans chaque sens. Sur la courbe de
l'\cref{ex:rc:commodity-and-energy-derivatives:calib}, le plan intrinsèque injecte
deux unités par mois d'avril à août et soutire en décembre, janvier et février
(\cref{fig:rc:commodity-and-energy-derivatives:curve})~; sa valeur est de 634\,447~USD. L'\omterm{def:rc:commodity-and-energy-derivatives:extrinsic}{intrinsèque glissant}
sur 4\,000 courbes simulées vaut 699\,550~USD, le Monte Carlo à moindres carrés
698\,182~USD~: une \omterm{def:rc:commodity-and-energy-derivatives:extrinsic}{valeur extrinsèque} d'environ 10~\% de l'intrinsèque à une
fréquence de décision mensuelle. Des décisions quotidiennes, que peut prendre une installation rapide, en ajoutent davantage.
\end{example}

\begin{omfigure}
\begin{tikzpicture}
\begin{axis}[width=11.5cm,height=5.4cm,
  xlabel={mois de livraison},ylabel={terme (\$/MMBtu)},
  axis y line*=left,
  xtick={1,2,3,4,5,6,7,8,9,10,11,12},xticklabels={avr.,mai,juin,juil.,août,sept.,oct.,nov.,déc.,janv.,févr.,mars},
  tick label style={font=\small},label style={font=\small},
  xmin=0.5,xmax=12.5,ymin=2.4,ymax=3.6,
  table/col sep=comma]
\addplot[omDef,very thick,mark=*,mark size=1.5pt] table[x=month,y=forward]{figdata/rates-credit-risk/16-commodity-and-energy-derivatives/curve.csv};\label{pl:rc:energy:fwd}
\end{axis}
\begin{axis}[width=11.5cm,height=5.4cm,
  axis y line*=right,axis x line=none,ylabel={inventaire (unités)},
  tick label style={font=\small},label style={font=\small},
  xmin=0.5,xmax=12.5,ymin=0,ymax=12,
  legend columns=2,legend style={font=\footnotesize,at={(0.5,-0.25)},anchor=north,draw=none,fill=none,/tikz/every even column/.append style={column sep=8pt}},
  table/col sep=comma]
\addlegendimage{/pgfplots/refstyle=pl:rc:energy:fwd}\addlegendentry{courbe à terme}
\addplot[omThm,very thick,const plot mark mid] table[x=month,y=inventory]{figdata/rates-credit-risk/16-commodity-and-energy-derivatives/curve.csv};
\addlegendentry{inventaire intrinsèque en fin de mois}
\end{axis}
\end{tikzpicture}

\omcaption{La courbe à terme illustrative (axe de gauche) et l'inventaire du plan intrinsèque
en fin de chaque mois (axe de droite)~: on remplit pendant les mois d'été bon marché, on vide pendant les
mois d'hiver chers. Données~: le tutoriel du chapitre.}\label{fig:rc:commodity-and-energy-derivatives:curve}
\end{omfigure}

\section{Contrats swing}

Un contrat swing (livre~3, chapitre~12) est un stockage sans injection~: le détenteur dispose d'un nombre de droits
de prendre du gaz à un prix fixe, limités par période et au total. Le
même programme dynamique le valorise, l'inventaire comptant les droits restants.

\begin{example}[Un swing d'hiver]\label{ex:rc:commodity-and-energy-derivatives:swing}
Droits de prendre jusqu'à 200\,000 MMBtu par mois de novembre à mars, soit 600\,000 au total, à
3.10~dollars~: sur la courbe à terme, le meilleur plan prend 200\,000 en décembre, janvier et
février, une valeur intrinsèque de 1.84 par unité de 100\,000 MMBtu~; le Monte Carlo à moindres carrés
donne 4.37. Dix calls mensuels, deux par mois sans la limite totale, valent 6.11, une
borne supérieure~: c'est la limite sur le total qui rend le swing moins cher qu'un strip.
\end{example}

\begin{dated}{2026-09}{Février 2021}\label{dat:rc:commodity-and-energy-derivatives:uri}
L'Energy Information Administration américaine a indiqué que, pendant la vague de froid de février
2021, le hub Waha a atteint 206.19~dollars par MMBtu le 16 février et SoCal Citygate 144.00
le 12 février, et qu'Oneok Gas Transportation, dans l'Oklahoma, a affiché en moyenne 1\,192~dollars le
17 février selon les données de NGI, contre 2.91 la première semaine du mois~; la
production de gaz a chuté à cause du gel des puits. Henry Hub a affiché en moyenne 5.35~dollars sur le mois, contre
2.71 en janvier.
\end{dated}

\begin{omfigure}
\begin{tikzpicture}
\begin{axis}[width=11.5cm,height=5.2cm,
  ylabel={Cours spot Henry Hub (\$/MMBtu)},
  xlabel={De décembre 2020 à avril 2021},
  xtick={-31,0,31,59,90,120},xticklabels={déc.,janv.,févr.,mars,avr.,mai},
  tick label style={font=\small},label style={font=\small},
  xmin=-31,xmax=120,ymin=0,ymax=26,grid=major,grid style={gray!25},
  table/col sep=comma]
\addplot[omDef,thick,mark=*,mark size=0.8pt] table[x=day,y=price]{figdata/rates-credit-risk/16-commodity-and-energy-derivatives/henryhub.csv};
\node[omThm,font=\footnotesize,anchor=west] at (axis cs:49,23.86) {23.86 le 17 février};
\end{axis}
\end{tikzpicture}

\omcaption{Prix spot quotidien du Henry Hub, de décembre 2020 à avril 2021. La vague de froid
a fait passer le prix de moins de 3~dollars à 23.86~dollars pendant une seule journée, avant de
le ramener sous 3~dollars en moins d'une semaine~: une valeur que seule une installation capable de livrer en quelques jours peut capturer. Données~: Energy Information
Administration américaine, série RNGWHHD (domaine public).}\label{fig:rc:commodity-and-energy-derivatives:hh}
\end{omfigure}

\section{Tutoriel~: stockage sous une courbe à deux facteurs}\label{tut:rc:commodity-and-energy-derivatives:tutorial}

\begin{tutorial}
\textbf{Objectif.} Calibrer le modèle à deux facteurs, pricer une option sur calendar spread et valoriser une
installation de stockage et un contrat swing. \textbf{État final~:} les résultats des
\cref{ex:rc:commodity-and-energy-derivatives:calib,ex:rc:commodity-and-energy-derivatives:cal,ex:rc:commodity-and-energy-derivatives:storage,ex:rc:commodity-and-energy-derivatives:swing}
et les quatre figures.
\begin{tutsteps}
\item \textbf{Calibrer}~: \texttt{model()} ajuste $\kappa$, $\sigma_s$, $\sigma_l$ à
\texttt{vol\_quotes()}.
\item \textbf{Spread}~: \texttt{calendar\_spread(0.5)}~; \texttt{kirk\_table()} par strike.
\item \textbf{Stockage}~: \texttt{storage()} donne les valeurs intrinsèque, \omterm{def:rc:commodity-and-energy-derivatives:extrinsic}{intrinsèque glissant} et
de Monte-Carlo par moindres carrés.
\item \textbf{Swing et storm}~: \texttt{swing()}, \texttt{storm()}~;
\texttt{fig\_rc\_energy.py} écrit les graphiques.
\end{tutsteps}
\textbf{Pour aller plus loin.} Doubler le taux de soutirage et observer quelle \omterm{def:rc:commodity-and-energy-derivatives:extrinsic}{valeur extrinsèque}
l'installation plus rapide ajoute~; fixer $\sigma_s$ à zéro et constater que la \omterm{def:rc:commodity-and-energy-derivatives:extrinsic}{valeur extrinsèque}
du stockage diminue.
\end{tutorial}

\section{Construction~: le modèle énergétique}\label{bld:rc:commodity-and-energy-derivatives:energymodel}

\begin{build}
\bfield{Rôle} Le moteur de matières premières de la société~: simulation de la courbe à terme pour les moteurs
de risque des parties~III et~IV, \omterm{def:rc:commodity-and-energy-derivatives:spread}{options sur spread} et valorisation de la
flexibilité physique.
\bfield{Interface} \texttt{TwoFactor(kappa, sigma\_s, sigma\_l, rho)} avec \texttt{var},
\texttt{cov}, \texttt{implied\_vol}, \texttt{simulate}, \texttt{forward}~;
\texttt{calibrate}~; \texttt{margrabe}, \texttt{kirk}, \texttt{spread\_mc}~;
\texttt{Facility}~; \texttt{intrinsic}, \texttt{lsm}, \texttt{rolling\_intrinsic},
\texttt{spot\_paths}~; \texttt{swing\_facility}.
\bfield{Règles} Une décision par période de livraison~; unités d'inventaire entières~; le spot d'une période
est son terme à la livraison~; simulation antithétique.
\bfield{Tests d'acceptation} \texttt{code/firm/energymodel/tests/}~: les termes simulés sont des
martingales dont la variance est celle du modèle~; la calibration retrouve les paramètres~; Kirk coïncide avec
Margrabe à strike nul et reste à moins de 1~\% de Monte-Carlo~; une courbe plate n'a pas de
valeur intrinsèque de stockage~; sans volatilité, la méthode de Monte-Carlo par moindres carrés égale la valeur intrinsèque~; l'\omterm{def:rc:commodity-and-energy-derivatives:extrinsic}{intrinsèque glissant}
et la méthode de Monte-Carlo par moindres carrés dépassent l'intrinsèque~; l'intrinsèque d'un swing est la somme de
ses meilleurs exercices.
\bfield{Pour aller plus loin} Décisions quotidiennes au sein du mois~; ratchets (taux dépendant de
l'inventaire)~; un facteur de sauts pour les pics~; la base spot-terme d'un hub physique.
\end{build}

\begin{omsources}
\item E.\ S.\ Schwartz, ``The stochastic behavior of commodity prices: implications for
valuation and hedging'', \emph{Journal of Finance} 52(3), 1997, 923--973.
\item E.\ Schwartz and J.\ E.\ Smith, ``Short-term variations and long-term dynamics in
commodity prices'', \emph{Management Science} 46(7), 2000, 893--911.
\item A.\ Boogert and C.\ de Jong, ``Gas storage valuation using a Monte Carlo method'',
\emph{Journal of Derivatives} 15(3), 2008, 81--98.
\item E.\ Kirk, ``Correlation in the energy markets'', dans R.\ Jameson (dir.), \emph{Managing Energy Price
Risk}, Risk Publications and Enron, London, 1995, 71--78.
\item US Energy Information Administration, \emph{Today in Energy}, ``Cold weather brings
near record-high natural gas spot prices'', février 2021~; série quotidienne des cours spot du Henry Hub.
\end{omsources}

\section{Exercices}

\begin{exercise}[$\star$]\label{exo:rc:commodity-and-energy-derivatives:1}
Avec les paramètres calibrés, quelle est la volatilité instantanée d'un terme à un an
de sa livraison, et pourquoi est-elle inférieure à celle du terme à la livraison~?
\end{exercise}

\begin{exercise}[$\star$]\label{exo:rc:commodity-and-energy-derivatives:2}
Pourquoi la \omterm{def:rc:commodity-and-energy-derivatives:intrinsic}{valeur intrinsèque du stockage} est-elle nulle sur une courbe à terme plate, et
pourquoi l'installation vaut-elle encore quelque chose~?
\end{exercise}

\begin{exercise}[$\star$]\label{exo:rc:commodity-and-energy-derivatives:3}
Une option d'achat sur calendar spread de strike 0.50 sur des termes à 3.50 et 2.74 expire demain.
Que vaut-elle~?
\end{exercise}

\begin{exercise}[$\star\star$]\label{exo:rc:commodity-and-energy-derivatives:4}
Pourquoi la corrélation entre deux mois de livraison compte-t-elle pour une \omterm{def:rc:commodity-and-energy-derivatives:spread}{option sur spread},
et dans quel sens~?
\end{exercise}

\begin{exercise}[$\star\star$]\label{exo:rc:commodity-and-energy-derivatives:5}
Pourquoi l'\omterm{def:rc:commodity-and-energy-derivatives:extrinsic}{intrinsèque glissant} ne peut-il jamais valoir moins que l'intrinsèque~?
\end{exercise}

\begin{exercise}[$\star\star$]\label{exo:rc:commodity-and-energy-derivatives:6}
Pourquoi un contrat swing vaut-il moins que le strip d'options d'achat mensuelles qu'il contient~?
\end{exercise}

\begin{exercise}[$\star\star\star$]\label{exo:rc:commodity-and-energy-derivatives:7}
\emph{Programmation.} Calculer l'erreur de Kirk aux strikes 0 et 0.5 et expliquer la différence.
\end{exercise}

\begin{exercise}[$\star\star\star$]\label{exo:rc:commodity-and-energy-derivatives:8}
\emph{Trouver la faille.} «~Notre portefeuille de stockage est valorisé à l'intrinsèque et entièrement couvert par des
forwards, donc février 2021 ne représentait pour nous ni un risque ni une opportunité.~»
\end{exercise}

\section{Problème~: La caverne de sel en février 2021}

\begin{problem}[{Problème du week-end --- une installation rapide et une installation lente pendant la vague de froid}]\label{pb:rc:commodity-and-energy-derivatives:1}
Deux installations détiennent chacune 500\,000 MMBtu près du Henry Hub le 8 février 2021. La caverne de sel
peut soutirer 100\,000 MMBtu par jour, le réservoir épuisé 10\,000. Entre le 8 et le 19 février,
chacune peut vendre, lors de ses meilleurs jours, au prix quotidien du Henry Hub, et rachète le même volume
au prix moyen de mars 2021 pour rétablir son plan~; le transfert de gaz coûte 2 cents par MMBtu dans chaque sens.

\medskip\noindent\textbf{Partie I --- Avant l'hiver.}
\begin{enumerate}
\item Donner les valeurs intrinsèque, intrinsèque glissante et de la méthode de Monte-Carlo
par moindres carrés de l'installation du chapitre.
\item Quelle est sa \omterm{def:rc:commodity-and-energy-derivatives:extrinsic}{valeur extrinsèque}, et d'où vient-elle~?
\item Quel paramètre du modèle de courbe influe le plus sur la \omterm{def:rc:commodity-and-energy-derivatives:extrinsic}{valeur extrinsèque}~?
\item Pourquoi un modèle mensuel sous-estime-t-il la \omterm{def:rc:commodity-and-energy-derivatives:extrinsic}{valeur extrinsèque} d'une caverne de sel~?
\item Comment couvrir la valeur intrinsèque au début de la saison~?
\end{enumerate}

\medskip\noindent\textbf{Partie II --- La tempête.}
\begin{enumerate}[resume]
\item Donner le cours du Henry Hub le premier février et à son sommet, ainsi que le jour du sommet.
\item Donner la moyenne de mars 2021 utilisée pour le rachat.
\item Donner le volume vendu par la caverne de sel, sur combien de jours, et son gain.
\item Donner la même chose pour le réservoir épuisé.
\item Pourquoi le rapport des gains n'est-il pas le rapport des débits de soutirage~?
\end{enumerate}

\medskip\noindent\textbf{Partie III --- Ce que voit un modèle.}
\begin{enumerate}[resume]
\item Combien d'écarts types séparent le cours de 3.76~dollars le 10 février de 23.86
le 17 février (quatre jours ouvrés) dans un modèle lognormal avec la volatilité court terme
de 69~\% du chapitre~?
\item Quelle caractéristique du modèle ferait apparaître de tels pics dans la simulation~?
\item Pourquoi la tempête a-t-elle fait monter la valeur des options de l'installation
mais non sa valeur intrinsèque pour la saison suivante~?
\item Quels contrats, en dehors du stockage, ont aussi gagné de l'argent~?
\item Qui était vendeur de cette flexibilité~?
\end{enumerate}

\medskip\noindent\textbf{Partie IV --- Jugement.}
\begin{enumerate}[resume]
\item Pourquoi les propriétaires de stockages rapides paient-ils cette option lorsque l'année moyenne
rapporte peu de \omterm{def:rc:commodity-and-energy-derivatives:extrinsic}{valeur extrinsèque}~?
\item Comment rendre compte du risque du portefeuille de stockage à un comité des risques~?
\item Que modifier dans le modèle de valorisation après février 2021~?
\item Énoncer le \emph{résultat nommé}~: les gains de l'installation rapide et de l'installation lente
pendant la tempête.
\item En une phrase~: que vaut une installation de stockage au-delà de son spread saisonnier~?
\end{enumerate}
\end{problem}

\section{Questions d'entretien}

\begin{interviewq}[$\star$ \iqroles{trader, chercheur}]\label{iq:rc:commodity-and-energy-derivatives:1}
Qu'est-ce que l'effet Samuelson, et comment un modèle à deux facteurs le produit-il~?
\end{interviewq}

\begin{interviewq}[$\star\star$ \iqroles{chercheur}]\label{iq:rc:commodity-and-energy-derivatives:2}
Comment valoriseriez-vous une \omterm{def:rc:commodity-and-energy-derivatives:spread}{option sur spread} à strike non nul~? Quelles approximations existent
et dans quels cas échouent-elles~?
\end{interviewq}

\begin{interviewq}[$\star\star$ \iqroles{trader}]\label{iq:rc:commodity-and-energy-derivatives:3}
Comment valorisez-vous un stockage de gaz~? Distinguer la valeur intrinsèque, l'\omterm{def:rc:commodity-and-energy-derivatives:extrinsic}{intrinsèque glissant}
et la pleine optionnalité.
\end{interviewq}

\begin{interviewq}[$\star\star$ \iqroles{développeur, chercheur}]\label{iq:rc:commodity-and-energy-derivatives:4}
Décrivez la méthode de Monte-Carlo par moindres carrés appliquée au stockage. Quels en sont les biais~?
\end{interviewq}

\begin{interviewq}[$\star\star\star$ \iqroles{risque}]\label{iq:rc:commodity-and-energy-derivatives:5}
Comment géreriez-vous le risque d'un portefeuille de stockage et de contrats swing pendant un hiver~?
\end{interviewq}

\begin{interviewq}[$\star\star\star$ \iqroles{chercheur}]\label{iq:rc:commodity-and-energy-derivatives:6}
Pourquoi un \omterm{def:rc:short-rate-models:g2}{modèle gaussien à deux facteurs} manque-t-il les pics de prix de l'électricité et du gaz
et qu'ajouteriez-vous~?
\end{interviewq}
